\documentclass[12pt,letterpaper]{article}
\usepackage[utf8]{inputenc}
\usepackage[T1]{fontenc}
\usepackage[spanish,es-noquoting]{babel}
\usepackage{mathptmx}
\usepackage[margin=2.5cm]{geometry}
\usepackage{setspace}
\usepackage{enumitem}
\usepackage{ragged2e}
\usepackage[hidelinks]{hyperref}
\onehalfspacing
\RaggedRight
\setlength{\parindent}{1.27cm}
\setlength{\parskip}{0pt}
\newenvironment{indice}{\begin{quote}\footnotesize\setstretch{1.1}\ttfamily}{\end{quote}}

\title{\vspace{-2.5cm}\large Avance para la revisión del 26 de septiembre de 2026}
\author{\normalsize Hesus García Cobos\\[2pt]
\footnotesize Maestría en Pedagogía, UPAEP\\
\footnotesize Director: Dr.\ Herminio Sánchez de la Barquera y Arroyo}
\date{}

\begin{document}
\maketitle
\thispagestyle{empty}

\section*{1. Nota sobre el cambio de marco}

Durante los meses anteriores construí el análisis comparado de las políticas de educación en
inteligencia artificial de siete países tomando como criterio de comparación seis valores de la
tradición confuciana. Ese criterio funcionó como instrumento de medición, pero llegué a una
objeción que no puedo sostener ante un jurado: elegí el conjunto de seis valores entre tres
candidatos por su poder de separación entre países, es decir por rendimiento estadístico y no por
pertinencia pedagógica. Un criterio de comparación elegido porque separa bien no es un criterio
de comparación justificado.

Cambio el criterio por las siete áreas de política que UNESCO enumera en la sección 6 de
\textit{AI and Education: Guidance for Policy-makers} (Miao, Holmes, Huang y Zhang, 2021). Son
una lista cerrada, numerada, específica de educación, publicada en español y escrita
explícitamente como plantilla para el diseño de política nacional. Usarla como rejilla de
comparación es el uso que sus autores previeron. La adopto además con dos anclajes: el Consenso
de Beijing (UNESCO, 2019), que da peso normativo a cada área, y los Marcos de Competencias de IA
de UNESCO (2024), que aportan el instrumento con el que valoro la competencia docente.

El cambio tiene un antecedente académico que me permite situar el trabajo en una conversación y
no en un vacío. Schiff (2022) analizó veinticuatro estrategias nacionales de inteligencia
artificial, cerca de mil quinientas páginas, con análisis temático de contenido, y propuso la
distinción entre \textit{educación para la inteligencia artificial} e \textit{inteligencia
artificial para la educación}. Mi tesis extiende ese diseño a siete países con documentos de 2020
a 2026 y con la rejilla de UNESCO como esquema de codificación, que es algo que no encontré hecho
en la literatura revisada.

Adopto la rejilla de UNESCO con reserva y lo digo en el marco conceptual. Mochizuki, Bruillard y
Bryan (2025) muestran que la orientación de UNESCO sobre inteligencia artificial en educación
descansa en una lógica tecno-solucionista y en una base de autores anglosajones y del sector
privado. Uso la rejilla como criterio de comparación, no como programa normativo que suscribo.

\section*{2. Índice del Capítulo 2}

\begin{indice}
2. Marco conceptual\\
\hspace*{1em}2.1 Antecedentes: estado del arte\\
\hspace*{2em}2.1.1 El análisis comparado de estrategias nacionales de IA\\
\hspace*{2em}2.1.2 La educación en las estrategias nacionales: el estudio de Schiff\\
\hspace*{2em}2.1.3 Rejillas de comparación aplicadas a políticas escolares de IA\\
\hspace*{2em}2.1.4 El análisis comparado en español y en América Latina\\
\hspace*{2em}2.1.5 La crítica a la orientación de los organismos internacionales\\
\hspace*{2em}2.1.6 Síntesis: coincidencias, diferencias, tendencias y contradicciones\\
\hspace*{1em}2.2 Marco teórico\\
\hspace*{2em}2.2.1 La política educativa como objeto de comparación\\
\hspace*{3em}a) El método comparativo de Bereday\\
\hspace*{3em}b) Los niveles de comparación de Bray y Thomas\\
\hspace*{3em}c) La política como texto, discurso y efecto\\
\hspace*{2em}2.2.2 Los marcos internacionales de IA y educación\\
\hspace*{3em}a) El Consenso de Beijing\\
\hspace*{3em}b) La Recomendación sobre la ética de la IA\\
\hspace*{3em}c) Las siete áreas de política de UNESCO como criterio de comparación\\
\hspace*{2em}2.2.3 La competencia docente en inteligencia artificial\\
\hspace*{3em}a) La alfabetización en inteligencia artificial\\
\hspace*{3em}b) El Marco de Competencias de IA para Docentes\\
\hspace*{2em}2.2.4 Revisión de estudios sobre formación docente en IA
\end{indice}

\section*{3. Esbozo del estado del arte (2.1)}

\textbf{2.1.1} Jobin, Ienca y Vayena (2019) revisaron ochenta y cuatro documentos de directrices
éticas de inteligencia artificial y encontraron convergencia en cinco principios: transparencia,
justicia y equidad, no maleficencia, responsabilidad y privacidad. Fjeld y sus colaboradores
(2020) trabajaron con treinta y seis documentos y organizaron ocho temas, entre ellos el control
humano de la tecnología y la responsabilidad profesional. Ambos estudios establecieron que las
políticas de inteligencia artificial convergen en el vocabulario y divergen en la
operacionalización. Ninguno de los dos se ocupó de la educación.

\textbf{2.1.2} Schiff (2022) llenó parte de ese hueco. Codificó veinticuatro estrategias
nacionales en inglés con un códebook preliminar y mostró que las estrategias tratan la educación
sobre todo como formación de fuerza de trabajo para la inteligencia artificial, y muy poco como
uso de la inteligencia artificial en los procesos de enseñanza. Su distinción entre educación
para la inteligencia artificial e inteligencia artificial para la educación organiza mi pregunta
segunda.

\textbf{2.1.3} Kundu y Bej (2025) compararon las políticas escolares de inteligencia artificial
de China, Singapur, Finlandia y Estados Unidos con nueve componentes operativos, entre los que
figuran la integración curricular y ética, la formación docente dinámica y el financiamiento
equitativo. Su trabajo demostró que una rejilla de componentes sirve para comparar políticas
escolares, aunque la construyeron los propios autores en lugar de tomarla de un instrumento
internacional.

\textbf{2.1.4} En español, Martínez Rojas, Suárez Guerrero y Martínez Bello analizaron de forma
comparada las políticas universitarias de regulación de la inteligencia artificial en Estados
Unidos y España en la \textit{Revista Española de Educación Comparada}, con cuatro dimensiones:
docencia y aprendizaje, evaluación, investigación y directrices éticas. Lara Martínez y Basave
Torres compararon la gobernanza de la inteligencia artificial en la educación superior de México
y Costa Rica con indicadores de 2019 a 2023. Los dos estudios trabajaron con dos casos; ninguno
llegó a siete países.

\textbf{2.1.5} Mochizuki, Bruillard y Bryan (2025) analizaron el discurso de la orientación de
UNESCO sobre inteligencia artificial en educación y la red de autores que cita. Concluyeron que
refleja una lógica de solución técnica y una base bibliográfica anglosajona y privada. Su crítica
no inutiliza la rejilla, pero obliga a declarar de dónde viene y con qué reservas la uso.

\textbf{2.1.6} Los cinco bloques coinciden en que la comparación de políticas de inteligencia
artificial es viable con rejillas de categorías, y difieren en el origen de la rejilla: unos la
derivan del propio corpus, otros la construyen como autores. Ninguno de los revisados toma una
lista de un organismo internacional y la aplica como esquema de codificación a un conjunto de
países. La tendencia dominante es la comparación de dos a cuatro casos. La contradicción más
visible está entre el uso creciente de los marcos de UNESCO como referencia y la crítica
documentada a su orientación.

\section*{4. Esbozo del marco teórico (2.2)}

\textbf{2.2.1} Comparo políticas educativas, así que necesito el aparato de la educación
comparada. Bereday (1964) propone cuatro etapas (descripción, interpretación, yuxtaposición y
comparación) que estructuran la lógica del análisis. Bray y Thomas (1995) cruzan niveles
geográficos, grupos demográficos y aspectos de la educación; su marco me obliga a explicitar que
trabajo en el nivel nacional, sin desagregar por grupo demográfico, y sobre el aspecto de
política de inteligencia artificial. Rizvi y Lingard (2010) distinguen la política como texto,
como discurso y como efecto; mi análisis opera en los dos primeros planos y no mide efectos de
implementación, lo cual delimita el alcance de lo que puedo afirmar.

\textbf{2.2.2} Los marcos internacionales fijan el vocabulario que las políticas nacionales
adoptan, adaptan o ignoran. El Consenso de Beijing (UNESCO, 2019) organiza sus recomendaciones en
ocho áreas, entre ellas la planeación de políticas, la inteligencia artificial para fortalecer la
enseñanza y a los docentes, el desarrollo de valores y competencias para la vida y el trabajo, y
el uso ético, transparente y auditable de los datos educativos. La Recomendación sobre la ética
de la inteligencia artificial (UNESCO, 2021) da el piso normativo. Las siete áreas de la sección
6 de la guía para responsables de política (Miao et al., 2021) son mi criterio de comparación:
visión sistémica y prioridades estratégicas; principio rector; planeación interdisciplinaria y
gobernanza intersectorial; políticas y regulación para un uso equitativo, inclusivo y ético;
planes maestros de gestión, enseñanza, aprendizaje y evaluación; pilotaje, monitoreo y
evaluación; y fomento de innovaciones locales.

\textbf{2.2.3} La intervención se dirige a docentes, así que necesito una definición operativa de
su competencia en inteligencia artificial. Long y Magerko (2020) y Ng y sus colaboradores (2021)
conceptualizan la alfabetización en inteligencia artificial; Touretzky y sus colaboradores (2019)
la traducen a ideas fundamentales para la educación básica. El Marco de Competencias de IA para
Docentes (UNESCO, 2024) organiza quince bloques en cinco aspectos por tres niveles de
progresión. Los aspectos son mentalidad centrada en el ser humano, ética de la inteligencia
artificial, fundamentos y aplicaciones, pedagogía de la inteligencia artificial, e inteligencia
artificial para el desarrollo profesional. De ese marco derivo el instrumento de autoevaluación
previa y posterior de mi intervención.

\section*{5. Pendientes que traigo a la revisión}

\begin{enumerate}[leftmargin=1.5em,itemsep=4pt]
\item Conseguir el artículo de Schiff (2022) para copiar las once etiquetas de su códebook. Está
tras muro de pago y no quiero reconstruirlas de memoria.
\item Confirmar con usted cómo se articula el análisis comparado con el apartado 1.2, ``Nuevo
contacto con la realidad''. Mi lectura es que el análisis documental de los siete países es la
fase de diagnóstico del primer ciclo de investigación-acción, y que las siete áreas de UNESCO son
las categorías que pide el apartado 1.3.
\item Definir la cohorte de docentes con la que aplicaré el plan de intervención. Es el punto con
el plazo más largo y por eso quiero cerrarlo primero.
\end{enumerate}

\section*{Referencias}

\begin{sloppypar}\setlength{\parindent}{0pt}\setlength{\parskip}{6pt}\small

Bereday, G. Z. F. (1964). \textit{Comparative method in education}. Holt, Rinehart and Winston.

Bray, M., y Thomas, R. M. (1995). Levels of comparison in educational studies. \textit{Harvard
Educational Review, 65}(3), 472--490.

Fjeld, J., Achten, N., Hilligoss, H., Nagy, A., y Srikumar, M. (2020). \textit{Principled
artificial intelligence: Mapping consensus in ethical and rights-based approaches to principles
for AI} (Berkman Klein Center Research Publication No. 2020-1). \url{https://ssrn.com/abstract=3518482}

Jobin, A., Ienca, M., y Vayena, E. (2019). The global landscape of AI ethics guidelines.
\textit{Nature Machine Intelligence, 1}(9), 389--399. \url{https://doi.org/10.1038/s42256-019-0088-2}

Kundu, A., y Bej, T. (2025). AI policies in school education: A comparative study on China,
Singapore, Finland, and the US. \textit{Journal of Science and Technology Policy Management}.
\url{https://doi.org/10.1108/JSTPM-06-2024-0218}

Lara Martínez, O. R., y Basave Torres, R. I. Educación superior y gobernanza de la inteligencia
artificial: Un análisis comparativo entre México y Costa Rica. \textit{RIDE, 17}(33).
\url{https://doi.org/10.23913/ride.v17i33.3107}

Long, D., y Magerko, B. (2020). What is AI literacy? Competencies and design considerations.
\textit{Proceedings of the 2020 CHI Conference on Human Factors in Computing Systems}.

Martínez Rojas, Á., Suárez Guerrero, C., y Martínez Bello, V. Regulación de la inteligencia
artificial en la Educación Superior: Un análisis comparativo exploratorio de las políticas
universitarias. \textit{Revista Española de Educación Comparada}, (50), 419--442.
\url{https://doi.org/10.5944/reec.50.2026.43461}

Miao, F., Holmes, W., Huang, R., y Zhang, H. (2021). \textit{AI and education: Guidance for
policy-makers}. UNESCO.

Mochizuki, Y., Bruillard, É., y Bryan, A. (2025). The ethics of AI or techno-solutionism?
UNESCO's policy guidance on AI in education. \textit{British Journal of Sociology of Education}.
\url{https://doi.org/10.1080/01425692.2025.2502808}

Ng, D. T. K., Leung, J. K. L., Chu, S. K. W., y Qiao, M. S. (2021). Conceptualizing AI literacy:
An exploratory review. \textit{Computers and Education: Artificial Intelligence, 2}, 100041.

Rizvi, F., y Lingard, B. (2010). \textit{Globalizing education policy}. Routledge.

Schiff, D. (2022). Education for AI, not AI for education: The role of education and ethics in
national AI policy strategies. \textit{International Journal of Artificial Intelligence in
Education, 32}(3), 527--563. \url{https://doi.org/10.1007/s40593-021-00270-2}

Touretzky, D., Gardner-McCune, C., Martin, F., y Seehorn, D. (2019). Envisioning AI for K-12:
What should every child know about AI? \textit{Proceedings of the AAAI Conference on Artificial
Intelligence, 33}(1), 9795--9799.

UNESCO. (2019). \textit{Beijing Consensus on artificial intelligence and education}.

UNESCO. (2021). \textit{Recomendación sobre la ética de la inteligencia artificial}.

UNESCO. (2024). \textit{Marco de competencias de la inteligencia artificial para docentes}.

\end{sloppypar}
\end{document}
